\abstracttext{%
China's Lunar and Planetary Data System (CLPDS) has accumulated open-access scientific datasets from seven deep space exploration missions, spanning Chang'e-1 through Chang'e-6 on the Moon and the Tianwen-1 Zhurong rover on Mars. However, extracting cross-mission planetary insights has historically been hindered by fragmented data formats, disparate processing levels (PDS3/PDS4), and domain-specific visualization hurdles. Here, we present the \textbf{CLPDS Planetary Exploration Data Visualization Suite}, a FAIR-compliant (Findable, Accessible, Interoperable, and Reusable) open-source framework and interactive WebGL visualization suite. Our platform integrates multi-scale planetary observations: (1) interactive 3D Moon and Mars globes displaying precision landing sites and rover traverse telemetry; (2) ground-penetrating radar (GPR) radargram analysis simulating Chang'e-4 Lunar Penetrating Radar (LPR, 500~MHz) and Zhurong Rover Subsurface Penetrating Radar (RoRAD, 25~MHz) echograms with dynamic dielectric permittivity ($\varepsilon_r$) depth scaling; (3) in-situ Vis-NIR and Laser-Induced Breakdown Spectroscopy (LIBS) spectral processing with automated continuum removal, resolving South Pole-Aitken (SPA) basin mantle-derived pyroxene/olivine absorption and Chang'e-5 in-situ hydroxyl ($OH/\text{H}_2\text{O}$) absorption at $2.85~\mu\text{m}$; and (4) Martian boundary layer meteorological (Mars Climate Station) and surface crustal magnetic field (RoMAG) diurnal dynamics in Utopia Planitia. This open platform bridges planetary geophysics, in-situ resource utilization (ISRU), and space life sciences, offering an accessible foundation for deep space exploration and bioregenerative habitat planning.
}

\keywordstext{CLPDS, Chang'e missions, Tianwen-1, Zhurong rover, Ground Penetrating Radar, Planetary Spectroscopy, Lunar Farside, FAIR Data}
